%%%%%%%%%%%%%%%%%%%%%%%%%%%% Document Setup %%%%%%%%%%%%%%%%%%
\documentclass[10pt]{article}
\usepackage[T2A,T1]{fontenc}
\usepackage[utf8]{inputenc}
\usepackage[russian,english]{babel}
\usepackage{CJKutf8}
\AtBeginDvi{\input{zhwinfonts}}
%\usepackage{polyglossia}
\makeatletter
\newlength{\bibhang}
\setlength{\bibhang}{1em}
\newlength{\bibsep}
 {\@listi \global\bibsep\itemsep \global\advance\bibsep by\parsep}
%\newenvironment{bibsection}%
 %       {\vspace{-\baselineskip}\begin{list}{}{%
   %    \setlength{\leftmargin}{\bibhang}%
     %  \setlength{\itemindent}{-\leftmargin}%
       %\setlength{\itemsep}{\bibsep}%
       %\setlength{\parsep}{\z@}%
        %\setlength{\partopsep}{0pt}%
        %\setlength{\topsep}{0pt}}}
        %{\end{list}\vspace{-.6\baselineskip}}
\makeatother
\reversemarginpar
\usepackage{geometry}\geometry{a4paper,total={170mm,257mm},left=1.18in,right=1.18in,top=0.79in,bottom=0.60in}
%\usepackage[paper=letterpaper,
           %\includefoot, % Uncomment to put page number above margin
   %         marginparwidth=0.6in, % Length of section titles
         %   marginparsep=0.05in, % Space between titles and text
      %      margin=1.18in, % 1 inch margins
            %includemp]{geometry}
%\usepackage{eurosans}
\setlength{\parindent}{0in}
\usepackage{enumerate} 
\usepackage{etaremune}
\usepackage{calc}
\usepackage{bbding}
\usepackage{paralist}
\usepackage{fancyhdr,lastpage}
\pagestyle{fancy}
%\pagestyle{empty} % Uncomment this to get rid of page numbers
\fancyhf{}\renewcommand{\headrulewidth}{0pt}
\fancyfootoffset{\marginparsep+\marginparwidth}
\newlength{\footpageshift}
\setlength{\footpageshift}
          {0.5\textwidth+1.\marginparsep+0.5\marginparwidth-2in}
\lfoot{\hspace{\footpageshift}%
       \parbox{4in}{\, \hfill %
                    \arabic{page} of \protect\pageref*{LastPage} % +LP
% \arabic{page} % -LP
                    \hfill \,}}
\usepackage{color,hyperref}
\usepackage{graphics}
\usepackage{wrapfig}
\usepackage{longtable}
\usepackage[usenames,dvipsnames]{xcolor}
\definecolor{darkblue}{rgb}{0.0,0.0,0.3}
\definecolor{mygreen}{RGB}{44,204,148}
\definecolor{darkgreen}{RGB}{13,107,27}
\definecolor{lightgreen}{RGB}{0,170,37}
\definecolor{provinces}{RGB}{11,100,132}
\definecolor{mypurple}{RGB}{77,24,137}
\definecolor{haiyan}{RGB}{8,150,178}
\definecolor{anticred}{RGB}{165,62,11}
\hypersetup{colorlinks,breaklinks,
            linkcolor=darkblue,urlcolor=darkblue,
            anchorcolor=darkblue,citecolor=darkblue}
%%%%%%%%%%%%%%%%%%%%%%%% End Document Setup %%%%%%%%%%%%%%%%%%%
%%%%%%%%%%%%%%%%%%%%%%%%%%% Helper Commands %%%%%%%%%

\newcommand{\makeheading}[1]%
        {\hspace*{-\marginparsep minus \marginparwidth}%
         \begin{minipage}[t]{\textwidth+\marginparwidth+\marginparsep}%
                {\large \bfseries #1}\\[-0.15\baselineskip]%
                 \rule{\columnwidth}{1pt}%
         \end{minipage}}
\renewcommand{\section}[2]%
        {\pagebreak[2]\vspace{1.3\baselineskip}%
         \phantomsection\addcontentsline{toc}{section}{#1}%
         \hspace{0in}%
         \marginpar{
         \raggedright \scshape #1}#2}
% An itemize-style list with lots of space between items
\newenvironment{outerlist}[1][\enskip\textbullet]%
        {\begin{itemize}[#1]}{\end{itemize}%
         \vspace{-.6\baselineskip}}
% An environment IDENTICAL to outerlist that has better pre-list spacing
% when used as the first thing in a \section
\newenvironment{lonelist}[1][\enskip\textbullet]%
        {\vspace{-\baselineskip}\begin{list}{#1}{%
        \setlength{\partopsep}{0pt}%
        \setlength{\topsep}{0pt}}}
        {\end{list}\vspace{-.6\baselineskip}}

% An itemize-style list with little space between items
\newenvironment{innerlist}[1][\enskip\textbullet]%
        {\begin{compactitem}[#1]}{\end{compactitem}}
% An environment IDENTICAL to innerlist that has better pre-list spacing
% when used as the first thing in a \section
\newenvironment{loneinnerlist}[1][\enskip\textbullet]%
        {\vspace{-\baselineskip}\begin{compactitem}[#1]}
        {\end{compactitem}\vspace{-.6\baselineskip}}
% To add some paragraph space between lines.
% This also tells LaTeX to preferably break a page on one of these gaps
% if there is a needed pagebreak nearby.
\newcommand{\blankline}{\quad\pagebreak[2]}
% Uses hyperref to link DOI
\newcommand\doilink[1]{\href{http://dx.doi.org/#1}{#1}}
\newcommand\doi[1]{doi:\doilink{#1}}
% For \url{SOME_URL}, links SOME_URL to the url SOME_URL
\providecommand*\url[1]{\href{#1}{#1}}
% Same as above, but pretty-prints SOME_URL in teletype fixed-width font
\renewcommand*\url[1]{\href{#1}{\texttt{#1}}}

\include{tikz-3dplot}
\usepackage{tikz,tikz-3dplot}
\usepackage{amssymb}
\usepackage{xifthen}
\usepackage{pstricks,pst-solides3d}
\usetikzlibrary{chains}
\usepackage{pgfplots}
\pgfplotsset{
compat=1.16,
    compat/path replacement=1.5.1,
}
\usepackage{pgfmath}
\usetikzlibrary{matrix}
\usetikzlibrary{arrows,decorations.pathmorphing,backgrounds,positioning,fit,petri}
\usepgfmodule{decorations}
\usepgflibrary{decorations.shapes} 
\usetikzlibrary[decorations.shapes] 
\usetikzlibrary{mindmap}
\usepgflibrary{patterns}
\usetikzlibrary[patterns]
\usetikzlibrary{petri}
\usepgflibrary{plothandlers}
\usetikzlibrary{plothandlers}
\usepackage{pgfpages}
\usepgflibrary{arrows}
\usetikzlibrary{arrows}
\usetikzlibrary{trees}
\usetikzlibrary{topaths}
\usetikzlibrary{3d}
\usetikzlibrary{positioning}
\usetikzlibrary{fit}
\usepgflibrary{shapes.misc}
\usetikzlibrary{shapes.misc}
\usepgflibrary{shadings}
\usetikzlibrary{shadings}
\usepgfmodule{shapes}
\usepgfplotslibrary{external}
\usetikzlibrary{pgfplots.external}
\usepgfplotslibrary{colorbrewer}
\usepgfplotslibrary{patchplots}
\usetikzlibrary{plotmarks}
\usepackage{color} 
\usepackage{xcolor} 
%\pgfpagesuselayout{resize to}[a4paper,border shrink=20mm]
\pgfplotsset{colormap={cool}{rgb255(0cm)=(255,255,255); rgb255(1cm)=(0,128,255); rgb255(2cm)=(255,0,255)}}
  \pgfplotsset{colormap={violet}{rgb255=(25,25,122) color=(white) rgb255=(238,140,238)}}
  \pgfplotsset{  colormap={bluered}{rgb255(0cm)=(0,0,180); rgb255(1cm)=(0,255,255); rgb255(2cm)=(100,255,0);rgb255(3cm)=(255,255,0); rgb255(4cm)=(255,0,0); rgb255(5cm)=(128,0,0)}}
 \pgfplotsset{/pgfplots/colormap={winter}{rgb255=(0,0,255) rgb255=(0,255,128)}}
 \pgfplotsset{/pgfplots/colormap={spring}{rgb255=(255,0,255) rgb255=(255,255,0)}}
 \pgfplotsset{colormap={CM}{rgb=(0,0,1) color=(black) rgb255=(238,140,238)}}
 
\usetikzlibrary{spy}
 \pgfplotsset{every patch/.style={miter limit=50},}
\usepgfplotslibrary{colormaps}


\begin{document} %  <-- begin document
\selectlanguage{russian}
\begin{tikzpicture} %  <-- begin 1st picture
	\begin{axis}[small,colorbar sampled line,title=Experiment №5802,xlabel={\tiny{Time, h}}, ylabel={\tiny{Immersion depth, mm}}, zlabel={\tiny{$C_{eq}$, MPa}}, xlabel style={sloped like x axis},ylabel style={sloped},colormap name=bluered,grid=major,view={210}{30},minor x tick num=10,minor y tick num=10] %  <-- plot annotations [in square brackets: axes titles, ticks, plot rotation]
  	  \addplot3+ [mesh,scatter, samples=48,domain=0:1, ] coordinates {
(0,0.25,0.38) (20,0.52,0.20) (40	,0.58,0.19)  (60,0.60,0.18) (80,0.61,0.17) (100,0.62,0.16) (120,0.63,0.15) (140,0.64,0.14)
(0.5,0.4,0.38) (1.5,0.45,0.20) (2.0,0.46,0.19) (2.5,0.47,0.18) (3.0,0.49,0.17) (3.5,0.50,0.16) (4,0.55,0.15) (5,0.60,0.14) 
(0,0.0,0.25) (10,0.15,0.23) (20,0.20,0.22) (30,0.30,0.21) (40,0.40,0.20) (50,0.50,0.19) (55,0.60,0.18) (60,0.70,0.17)
(0,0.0,0.35) (2.0,0.04	,0.31) (4,0.08,0.29) (5,0.10,0.27) (6,0.12,0.25) (7,0.14,0.23) (8,0.15,0.21) (10,0.16,0.20)
(0,0.0,0.34) (2.0,0.10,0.31) (3.0,0.20,0.28) (4,0.30,0.26) (5,0.35	,0.24) (6,0.4,0.23) (7,0.45,0.22) (8,0.50,0.21)
(0,0.0,0.33) (5,0.20,0.30) (7,0.30,0.29) (11	,0.35,0.28) (13,0.40,0.26) (15,0.45,0.24) (17,0.47,0.22) (19,0.52,0.20)  
};
	\end{axis}
\end{tikzpicture} %  <-- end 1st picture
	\hskip 5pt %  <-- add space between the subplots (between 1st and 2nd figures)	
\begin{tikzpicture}  %  <--  begin 2nd picture
	\begin{axis}[small=4cm,colorbar sampled line,title=Experiment №5803,xlabel={\tiny{Time, h}},  %  <-- annotate axes n the plot, adjust font size: make text tiny. 
    ylabel={\tiny{Immersion depth, mm}}, zlabel={\tiny{($C_{eq}$), MPa}}, xlabel style={sloped like x axis}, %  <-- here we annotate axes in parallel to the axes rotation, i.e. make the text parallel to the axe
    ylabel style={sloped like y axis}, zticklabel style={rotate=90, %  <-- here we rotate y axe label  to 90 degree.
    	/pgf/number format/precision=3, %  <-- get rid of the degrees instead of many zeroes after comma, like that: 10*(-4)
    	/pgf/number format/fixed, %  <-- fix numeric format (w/o floating point)
        /pgf/number format/fixed zerofill, %  <-- here we annotate all zeroes, even in very small numbers, like 0.000001
},yticklabel style={sloped like y axis},colormap/PuOr,grid=major,view={210}{30},minor x tick num=10,minor y tick num=10] %  <-- environment of the 2nd plot and [its description] in square brackets
 	   \addplot3+ [mesh,scatter, samples=42,domain=0:1, ] coordinates {
(0,0.6,0.110) (20,0.9,0.070) (40,1.0,0.062) (60,1.2,0.060) (80,1.3,0.059) (100,1.4,0.058) (120,1.5,0.057) %  <-- this is a table of sampling data in 3D spherical coordinates. Here is example of the compressive strength measurements of soil samples.
(0,0.50,0.110) (2,0.65,0.105) (4,0.70,0.100) (6,0.73,0.080) (8,0.77,0.075) (12,0.80,0.073) (16,0.85,0.070)
(0,0.40,0.110) (2,0.52,0.100) (4,0.62,0.090) (6,0.65,0.088) (8,0.70,0.085) (12,0.73,0.083) (16,0.75,0.080) 
(0,0.50,0.130) (2,0.58,0.110) (3,0.63,0.106) (5,0.65,0.103) (7,0.70,0.100) (10,0.72,0.095) (12,0.75,0.090) 
(0,0.40,0.170) (2,0.50,0.155) (4,0.55,0.150) (6,0.58,0.135) (8,0.62,0.120) (12,0.67,0.110) (16,0.70,0.100) 
(0,0.51,0.133) (2,0.59,0.115) (3,0.64,0.104) (5,0.67,0.096) (7,0.72,0.090) (10,0.74,0.083) (12,0.77,0.080)
};
	\end{axis} %  <-- end plot environment
\end{tikzpicture} %  <-- end 2nd picture (on the right)
	\begin{flushleft}
		\textbf{Time series analysis of the variations of the long-term equivalent cohesion of soils ($C_{eq}$) under increasing depth of the samples. Series-II (№1-8). 3D: \LaTeX{}}
	\end{flushleft}
\end{document} %  <-- end document